\section{Contribuição}
\label{sec:contribuicao}

Para o sistema Banco de Especialistas, a principal contribuição foi a formalização 
do processo de solicitação, revisão e reversão das alterações cadastrais e de 
categoria. O módulo substituiu um fluxo baseado em e-mail e execução manual de 
SQL por um fluxo em que a intenção do pedido, a decisão autorizada e o efeito sobre o 
cadastro são registrados de forma consistente, com histórico completo e controle de 
acesso aplicado no servidor, conforme as Seções~\ref{sec:modulo-solicitacoes} e \ref{sec:snapshot-auditoria}. Como consequência, as operações 
passaram a ser auditáveis nas três dimensões — temporal, autoral e 
semântica, conforme a Seção~\ref{sec:snapshot-auditoria} —, e a reversão deixou de ser uma exclusão de histórico 
para se tornar um evento registrado, também na Seção~\ref{sec:snapshot-auditoria}.

Para a equipe técnica, o trabalho deixou três legados principais. O primeiro 
é uma base de testes reutilizável, inexistente antes do módulo, que estabeleceu uma 
linha de base de regressão para as funcionalidades existentes e pode ser estendida a 
desenvolvimentos futuros, conforme a Seção~\ref{sec:testes-validacao}. O segundo é um pipeline de integração contínua 
institucional, executado no Runner interno do Polo de Inovação, com estágios de lint 
e test e bloqueio de merge em caso de falha, conforme a Seção~\ref{sec:testes-validacao}. O terceiro é um padrão de 
separação de responsabilidades — controladores finos, Service Layer e validação na 
fronteira — que serve como referência para a evolução de outros módulos do sistema, 
em contraste com o padrão concentrado dos controladores legados, como descrito nas Seções~\ref{sec:arquitetura} e \ref{sec:modulo-solicitacoes}.

Para o processo de desenvolvimento, o trabalho constituiu uma demonstração prática da 
estratégia de Additive Change em um sistema legado, descrita na Seção~\ref{sec:additive-change}. A incorporação do módulo 
como um componente isolado, por meio de Blueprints e de extensões às rotas e aos stores 
existentes, preservou o comportamento das funcionalidades já consolidadas e reduziu o 
risco de regressão, evidenciando que é possível evoluir um sistema em uso sem reescrever suas estruturas centrais.

Para a formação do estagiário, o estágio proporcionou a consolidação de competências 
full-stack, ao percorrer todas as etapas do desenvolvimento — levantamento de requisitos, 
modelagem, implementação de back-end e front-end, testes e integração contínua — em um 
sistema real, com regras de negócio definidas e restrições institucionais. O contato 
com testes automatizados e com o pipeline de CI, em particular, ampliou a compreensão 
sobre a importância da verificação contínua e da rastreabilidade no desenvolvimento de software.